
\texttt{top(...)} decides where that split falls, so it is worth settling rather than
asserting. An earlier draft of this paragraph called it ``close to PromQL's
\texttt{topk}'', and that is false: Datadog ranks series by an aggregation over
\emph{all} values in the queried range and returns those series' whole trajectories,
while \texttt{topk} ranks independently at each evaluation instant, so a series can be in
one and not the other. They are different operators, and $139$ occurrences moving into
the inexpressible column would invert the comparison this section ends with.

They do not move, and the reason is instructive. PromQL \emph{can} denote the ranking,
using the \texttt{@} modifier to pin it to a single evaluation so that it does not drift
across the graph:
\begin{center}\footnotesize
\texttt{q and on($\overline{k}$) topk($n$, max\_over\_time(q[$R$] @ $T$))}
\end{center}
Read off a real Prometheus~3.15.0, this reproduces the selection that Datadog's
documented rule predicts --- aggregate each series over the whole range, take the top
$n$ --- and holds it constant across the range, while the naive \texttt{topk($n$, q)}
changes membership within the same workload at $n=2$
(\texttt{scripts\slash verify\_topk.py}; cached and regression-tested).
All $20$ cases agree: five ranking
functions, at $n=2$ and $n=3$, in both directions. The directions were the omission worth
finding. This corpus ranks ascending in $5$ of its $139$
occurrences, \texttt{topk} is the descending half of \texttt{top()}, and until
\texttt{bottomk} was asked for as well the measurement covered $134$ occurrences while
the sentence reporting it claimed $139$. The functions run the other way: the probe
covers five and the corpus uses four --- \texttt{mean} in $111$ of the calls,
\texttt{last} in $34$, \texttt{max} in $15$ and \texttt{sum} in $3$ --- so \texttt{min}
is probed because the candidate has to hold for it, not because a dashboard asks.

The Prometheus reading is one side of a comparison, and the other side of it was a
reading of Datadog's documentation --- the weaker kind of evidence, and the kind this
paper has twice found wrong, since the aggregation order and the bucket closure were both
read off the documentation and both wrong. So the vendor was asked the same question. The five series
go to a real Datadog as a gauge and \texttt{top()} itself is asked for its selection
(\texttt{scripts\slash datadog\_top.py}; cached and regression-tested), over the one
workload definition both scripts import so that the two halves cannot drift apart. All
$22$ selections are the ones the documented rule predicts, and two things the
documentation asserts rather than demonstrates are now readings as well. The selected
series come back over the \emph{whole} window and not only where they are winning: the
series that spikes once and then sits at $1$ returns all $20$ buckets, which is precisely
what a \texttt{topk}-shaped operator would not do. And the selection is relative to the
queried range --- over all twenty buckets the top two by \texttt{max} are the riser and
the spike, over the last ten the spike is outside the window and the answer is the riser
and the flat high series, while \texttt{mean}'s selection does not move under the same
narrowing. So \texttt{top(...)} is unimplemented work, and the catalogue's split stands.

What the expression costs is the interesting part, and that last reading is what prices
it. $R$ is the dashboard's time range and $T$ its end, and a Datadog query carries
neither --- the same omission as the rollup interval of \S\ref{sec:caveats}. That would
be a detail if the vendor's answer were insensitive to the range, and it is not: a
translator told the wrong range does not draw a slightly wrong graph, it selects
different series. An implementation of \texttt{top(...)} would therefore
reach the ceiling below only for a translator that is told the time range, and would
carry two more preconditions of exactly the kind \S\ref{sec:discussion} argues are the
real switching cost.

\texttt{PHI\_QUANTILE} ($34$) is the refusal this corpus comes closest to raising as an
information barrier in the main direction, and \S\ref{sec:ingestion} proves the
underlying fact rather than asserting it: a Datadog distribution metric answers
$\mathtt{p}NN$ from a sketch of the raw observations, Prometheus keeps those observations
bucketed with their sum and count, and no function of those is the quantile
(Corollary~\ref{cor:quantile}). The point is not that one system is more precise --- a
sketch is lossy too --- but that the two keep different lossy summaries of one stream and
neither recovers the other.

What this remark called ``one of the two information barriers this corpus raises'' it is
not, and the correction is \S\ref{sec:gaps}'s: the corollary's hypothesis is a
\emph{distribution} metric, and of the $34$ occurrences $20$ name metrics the vendor's
repository declares as something else, so the barrier does not apply to them. The
remaining $14$ name metrics the repository does not declare at all, which leaves them
undecided rather than blocked. So the sound reading of this refusal is a proved fact
about sketches and bucketings, a construct that fires on \texttt{p}$NN$ without testing
the fact's hypothesis, and $14$ occurrences where the hypothesis might hold and nothing
here says whether it does. PromQL's \texttt{histogram\_quantile} is not
the missing operator; it is a different approximation.

